%subfiles使用時はdocument環境内にサブファイルも記入する
%サブファイル内のプリアンブルはmain.texのビルド時は無視される
\documentclass[../../main.tex]{subfiles} %[]内にこのドキュメントからmain.texまでのpathを記入

\graphicspath{{../../fig/}} %このドキュメントから画像ディレクトリまでのpathを記述
\setcounter{chapter}{0} %チャプター番号の調整

%***************************************** ＜Main body＞ ********************************************%
\begin{document}

\setcounter{page}{0} %ページ番号をゼロに設定
\pagenumbering{arabic} %本文用のページ番号付け

\chapter{関連研究}\label{chap:introduction}

\section{Related Work}

Robots are expected to support human life and activities in a wide range of 
fields such as medicine, manufacturing, space exploration, and agriculture.  
However, many challenges remain for 
humans and robots to work collaboratively in the same space.  
One of these is the lack of mutual understanding of intentions and actions.  
Humans often find it difficult to accurately predict the capabilities and behaviors of robots,
while robots cannot fully interpret humans’ high-level goals.  
As a result, “gaps of execution” and “gaps of evaluation” arise, hindering smooth collaboration.  
Conventional interfaces have mainly relied on joysticks or keyboards, but visualization on 2D displays 
provides insufficient depth information, 
limiting the operator’s spatial awareness and sense of immersion.

To address these challenges, research on Human-Robot Collaboration (HRC) using Mixed Reality (MR) 
has made significant progress in recent years.  
By overlaying information about robot states and environmental conditions in MR space, 
users can perform more intuitive operations, 
and improvements in efficiency for teleoperation and collaborative tasks can be expected \cite{ref1,ref2}.  
In addition, many studies have explored applications of AR and VR.  
In VR, methods have been proposed to support skill acquisition by robots through 
imitation learning and motion data collection \cite{refX},  
while in AR, attempts have been made to assist workers’ understanding by presenting 
task procedures and robot states \cite{refY}.  
Nevertheless, these approaches still face limitations in interaction with the real 
environment and in providing a sense of immersion.  
In this context, MR, which can integrate the real world with virtual information, 
is receiving particular attention as a promising foundation technology for HRC.

Previous MR-based HRC research can largely be classified into two directions.

First, research focusing on \textbf{visualization of robot states}.  
This direction projects sensor data and robot internal states into MR space, 
thereby supporting users’ situational awareness and decision-making \cite{ref4}.  
While this approach contributes to enhancing environmental understanding and reducing cognitive load, 
direct intervention in robot motion planning and control remains limited.

Second, research focusing on \textbf{intuitive mapping of user operations}.  
By transforming hand tracking and gesture input into robot motions, these approaches enhance 
operability and immersion.

In this category, the work of Esaki et al. is representative.  
They proposed the \textit{Immersive Robot Teleoperation Based on User Gestures in Mixed Reality Space (IRT-MRO)} \cite{ref5}.  
In IRT-MRO, objects are placed in MR space, and users directly manipulate their positions and orientations.  
The manipulated results are transferred to the robot side, enabling the robot to 
execute the same operations in the physical world.  
This demonstrated the effectiveness of MR in improving the intuitiveness of robot teleoperation.  

However, IRT-MRO also leaves several issues unresolved.  
First, because the system simply transfers the object manipulation performed by the user in MR space, 
it does not take into account whether the robot can actually reach the target object or grasp it stably.  

Second, the posture of the robot arm is not always guaranteed to remain safe, 
and problems such as singularity or joint limit violations can occur.  

Third, when an object is located outside the reachable area of the robot, 
the system does not provide an appropriate relocation strategy for the robot base, 
leaving the operation incomplete.  

In summary, previous studies have clarified the potential of MR as an intuitive interface for HRC.  
Nevertheless, they still face challenges regarding (1) robust target inference, 
(2) ensuring robot posture and grasp stability, and (3) achieving autonomous yet 
interactive base relocation strategies.  
These issues become bottlenecks in realizing seamless human-robot cooperation in dynamic and complex environments.  

To address these challenges, this paper proposes a new framework called \textbf{Human–Robot Interaction with Correction and Suggestion (HRI-CAS)}.  
Unlike previous MR-based systems, HRI-CAS integrates:  
\begin{enumerate}
  \item \textbf{Target Bottle Inference (TBI)}: a module that infers the user’s 
  intended target based on multimodal features such as gaze, 
  head direction, and hand motion, with adaptive information-level control.
  \item \textbf{Base Relocation based on Inverse Reachability Map (BRI)}: 
  a module that determines optimal robot base positions when the target is unreachable, 
  by leveraging IRM-based optimization and MR-based user corrections.
  \item \textbf{Human-Intent Adaptive Grasp Manipulation (HIGM)}: a module that integrates Pregrasp design, nullspace optimization, 
  and finite state machines to adaptively align user intention with the robot’s kinematic constraints.
\end{enumerate}

Through these three modules, HRI-CAS enhances the robustness of target selection, 
guarantees stable robot motion, and provides a cooperative relocation mechanism, 
thereby advancing MR-based HRC beyond the limitations of IRT-MRO.



\end{document}